\titleformat{\section}{\Large\bfseries\filcenter}{}{0pt}{}

\addchaptertocentry{Ерөнхий дүгнэлт}

\section*{Ерөнхий дүгнэлт}
\markboth{Ерөнхий дүгнэлт}{Ерөнхий дүгнэлт}

Энэхүү ажлаар монгол хэлээр асуусан асуултыг банк бус санхүүгийн байгууллагын PostgreSQL өгөгдлийн сан дээр
гүйцэтгэх SQL болгон хөрвүүлэх системийг судалгааны үр дүнд тулгуурлан бүтээж, хяналттай туршилтаар үнэлэх бүрэн
дамжлагыг хэрэгжүүлэв.

\textbf{Онолын хувьд} (Бүлэг~\ref{Chapter1}) LLM-д суурилсан Text2SQL-ийн архитектур, хайлтаар баяжуулсан үүсгэлт,
rejection sampling ба GRPO-д суурилсан fine-tuning, agentic арга, бага нөөцтэй хэлний стратеги, үнэлгээний хэмжүүрийг
тоймлов. Монгол хэлний гол бэрхшээлийг тодорхойлов: залгамал үгзүй, кирилл ба латин галигийн холимог, монголоор
хадгалагдсан утга, загварын англи төвтэй байдал. Бэрхшээл бүрд тохирох хариу аргыг сонгов.

\textbf{Хэрэгжүүлэлтийн хувьд} (Бүлэг~\ref{Chapter2}) дараах бүрэлдэхүүнүүдийг бүтээв:
\begin{itemize}
  \item регистр, утас, хаяггүй 28 view бүхий text2sql давхарга, мөрийн түвшний хамгаалалт (RLS);
  \item 425 баганын монгол тайлбар, 45 нэр томьёо, 20 KPI бүхий мэдлэгийн давхарга;
  \item кирилл ба галигийн түлхүүртэй утгын хайлт;
  \item хоёр загварын өөрийн протокол, давхаргатай SQL хамгаалалт;
  \item засвар, санал хураалт, татгалзалт, tool ба шүүмжлэгчтэй agentic давталт;
  \item орчуулга, SFT ба GRPO-ийн дамжлага.
\end{itemize}
Бүрэлдэхүүн бүрийг нэг тохиргооны талбараар асааж унтраадаг тул нөлөөг нь тусад нь хэмжих боломжтой.

\textbf{Туршилтын хувьд} (Бүлэг~\ref{Chapter3}) 8 домэйн, 500 асуулттай, кирилл, латин галиг, англи, холимог
хувилбартай монгол eval set-ийг LLM агентуудаар бэлтгэж, бие даасан LLM хянагчаар шалгав. Үүнд хариулах боломжгүй 62, хоёрдмол 63 асуулт орсон. Gold бүрийг 9 өгөгдлийн
сан дээр баталгаажуулав. 50 тохиргоог (A1--G6) «нэг удаад нэг өөрчлөлт» зарчмаар төлөвлөж, EX, TS-EX, TrustSQL,
McNemar-аар харьцуулах дахин давтагдах дамжлага бүрдүүлэв. Үнэлгээний дамжлагын зөв байдлыг дараах аргуудаар
баталгаажуулав: хоёр загвар, хоёр өгөгдлийн сан дээрх oracle run (EX 100\%, ББСБ дээр TrustSQL 1.0), аюулгүй байдлын 84
халдлага (бүгд хаагдсан), 63 тест, кодын хоёр удаагийн бие даасан хяналт (68 олдвор; засагдсан эсвэл баримтжуулсан)
ба тайлангийн бие даасан хяналт. Мөн өмнөх ажлын 150 кейсийг шилжүүлж, хоёр
системийг ижил өгөгдөл, загвар, шалгуураар харьцуулах орчин бүрдүүлэв.

\ifnum\resultsavailable=1
\textbf{Үр дүнгийн хувьд} судалгааны асуултуудын хариултыг Бүлэг~\ref{Chapter3}-ийн хүснэгтүүдэд тулгуурлан энд
нэгтгэнэ. % end5896 — серверийн үр дүнгээс RQ1--RQ5-ын хариултыг бичих
\else
\textbf{Үр дүн.} Загварын бодит оноог GPU сервер дээр гаргана. Судалгааны асуултуудын (RQ1--RQ5) хариултыг
серверийн туршилтын дараа энд нэгтгэнэ. % end5896 — серверийн үр дүнгээс RQ1--RQ5-ын хариултыг бичих
\fi

\textbf{Цаашдын ажил.}
\begin{enumerate}
  \item Eval set-ийг ББСБ-ийн шинжээчээр хянуулж (зааврын «gold-ийг хоёр хүн шалгана» шаардлага), дараа нь серверийн
        туршилтыг ажиллуулах.
  \item Shadow горимд бодит хэрэглэгчийн асуулт цуглуулж, eval set-ийг бодит тархалтаар баяжуулах.
  \item Хэрэглэгчийн засвараас галиг хувиргагч сургах (F3), идэвхтэй тэмдэглэгээг (F9) бодит лог дээр ажиллуулах.
  \item Олон ээлжит tool орчинтой RL-ээр agentic загварыг (sqrl) монгол өгөгдөл дээр сургах.
  \item Монгол хэлэнд хүчтэй загварыг шүүмжлэгч ба орчуулагчаар хослуулах.
\end{enumerate}
